\section{讨论}
\label{sec:discussion}

\subsection{平台实际发现了什么}

在本案例中，平台得到的不是"更高的预测精度"，而是一组经过分层筛选的因素证据：序列层面指向 PAM 邻近区域（尤其第 18 位）作为候选序列相关因素；模型层面显示更宽的局部感受野带来稳定但有限的预测改善，并改变归因分布；环境层面则得到一个小效应且不稳健的结果；细胞背景层面显示同一因素的效应方向与幅度依赖细胞系。这些结论共同构成一条从"预测"到"候选因素"的证据链，而其中最重要的能力是：平台在统一框架内**同时**输出了稳定证据、模型特异性现象与不确定结果。

\subsection{为什么多模型证据比单模型更可信}

五类模型的归纳偏置不同：线性模型给出方向明确的主效应，树模型刻画非线性与交互，MLP 表示一般非线性组合，CNN 建模局部序列模式，Transformer 建模跨位置关系。本案例中，PAM 邻近区域的最高归因在四个模型类中重复出现，而线性模型的峰值位置与注意力谱的峰值位置不同。若只报告单一模型，这种差异不可见，读者会把某一种归因方式误当作事实。平台的做法是：把跨模型重复出现的结构作为较强候选，把仅在某类模型出现的现象标记为 model-dependent evidence。这不仅提高了结论的可信度，也明确了结论的适用范围。

\subsection{环境效应为什么弱，以及这一结果的价值}

表观遗传环境在本案例中提供的增量预测价值很小（$|\Delta R^{2}|<0.01$），边级置信区间多数跨 0，区组析因方差分析也未检出显著主效应（$p=0.056$--$0.663$）。可能的解释有三：其一，四条轨道以逐位点二值形式编码，信息容量有限；其二，环境信息与序列特征存在冗余，sequence-only 模型已能捕获相关信息；其三，不同模型类对环境通道的利用方式不同，跨模型平均会削弱单一模型的效应。三种口径（置换检验、区组 ANOVA、边级 bootstrap）结论不一致本身也说明该效应对依赖结构假设敏感。**这一结果的价值在于平台没有把它包装成阳性结论**：科学发现平台的可信度取决于它能否识别并如实报告"信息增量不足"的情形。

\subsection{第 18 位结果意味着什么}

第 18 位的证据由三部分构成：多个模型的位置归因（3/5 模型 top-3，跨模型平均谱第 1 位）、细胞系依赖的实测效率对比（C 相对 A：HCT116 $+0.090$、HeLa $+0.092$、HL60 $+0.028$、HEK293T $-0.015$）以及 CNN 替换敏感性（HCT116/HeLa 中 A/C 的替换效应明显高于 G/T，HEK293T 中四种碱基接近）。这提示一个候选序列相关模式，但需要严格区分两类解释：**位置效应**（该位置本身的序列语法重要性）与**碱基特异性效应**（该位置上特定核苷酸的效应）。本数据无法把二者完全分离，且 HEK293T 的反向可能部分来自其碱基组成差异（该细胞系第 18 位含 24.7\% T）。因此本文的结论是"具有上下文依赖性的候选序列关联"，而非"第 18 位 C 是普适决定因素"。

\subsection{感受野消融告诉我们什么}

配对实验显示更宽的卷积核稳定改善预测（$k=7-k=3$：$+0.056$，95\% CI 0.051--0.060；$k=5-k=3$：$+0.042$；$k=7-k=5$：$+0.013$），并使归因更集中于 PAM 邻近区域。这支持"更宽的局部序列上下文包含额外预测信息"的模型层面解释。但候选模式长度并未随核大小增加（三个配置的模式长度中位数均为 4\,nt），说明卷积核描述的是模型感受野，而不是生物学模式长度。这一区分具有普遍方法学意义：在把深度模型的滤波器或窗口长度解释为生物学单元之前，必须用独立于模型结构的证据加以检验。

\subsection{为什么细胞背景必须被建模而不是被平均}

同一环境因子在不同细胞系中方向与幅度不同（CTCF 在 HL60 反向；H3K4me3 在 HEK293T 与 HL60 方向相反），序列特征同样如此（第 18 位 C 的效应在 HEK293T 反转）。若把细胞系简单合并，会得到接近零的平均效应，从而掩盖上下文依赖结构。平台因此输出 context 标签（一致/依赖/冲突/不确定）而非单一平均值。同时必须克制：细胞系差异也可能来自数据量、标签分布或序列组成差异，因此本文只作"上下文依赖关联"的表述，不作机制断言。

\subsection{哪些结论仍未经验证}

本文的所有结果都属于计算证据：预测性能来自留出测试集，归因来自模型内部，统计检验来自对既有结果的重新分析，候选模式来自归因聚合与富集检验。**尚未完成**的工作包括：直接序列扰动实验、编辑效率的湿实验验证、独立生物学验证，以及在本批次中退化的真正 LOCO 重训。因此平台输出的是可检验假设，而不是已确认的生物学机制。

\subsection{如何由此设计后续实验}

候选优先级（表~\ref{tab:candidates}）给出了两条最小扰动路径：C1 以单碱基替换检验第 18 位的 C/A 差异（同一 sgRNA 其余序列保持不变，可同时在不同细胞系中测试以检验上下文依赖）；C2/C3 通过破坏 PAM 邻近的 4\,nt 候选核心检验模式层面的假设。两条路径均不依赖复杂的环境操控，便于在常规编辑效率测定体系中实施。相关的实验设计与预期读出方式见 \S\ref{sec:future}。

\subsection{面向赛道定位的能力对照}

从"AI 基因编辑与核酸工具设计"的角度，本工作的重点可以概括为三点：其一，**方法层面**使用异构模型族与多源特征（序列 + 表观通道 + 细胞背景），并把每个分析单元的可解释输出统一到可追溯的结果文件；其二，**发现层面**给出可排序的候选因素与候选序列模式，而不是只给出预测分数；其三，**验证层面**明确区分已完成的计算验证与尚未开展的湿实验验证，并给出可执行的最小扰动方案。平台当前实现的是候选优先级排序（candidate prioritization / ranking），不是序列从头生成；这一点在方法（\S\ref{sec:methods}）与候选章节中均如实说明。
